% 来源：OEBPS/Text/Section0001_split_008.xhtml#sigil_toc_id_60

\BookChapter{46}
\section[第46章\q 最终的秘密]{第46章\\ 最终的秘密}

\begin{BookQuote}
我们靠我们得到的东西生存，我们靠我们给予的东西生活。

\hfill ——温斯顿·丘吉尔
\end{BookQuote}

现在我们要一起迈向这次投资之旅的最后一步，我想请你仔细考虑一下：在这个世界上，最能让你激情燃烧的是什么事情？你最在乎的是什么东西？什么事情让你兴奋激动？留给后人什么遗产会让你两眼放光？今天做什么事会让你感到骄傲自豪？你采取什么样的行动，会给个人精神发信号，表明你的人生过得很好？如果你真的受到了启发，你愿意创造什么，或者给予什么？

所有这些问题都会带我们更加接近真正财富的最终秘密。但是这里有一笔交易，其中一部分关键因素也许看起来完全违背直觉。我们花了大量时间来谈论如何掌握金钱、储蓄、投资，打造达到关键数量的财富规模，从而能够最终创造出自由，提升你的生活品质，但是与此同时我们受到的教导是，钱不能买来幸福。正如一项研究想要证实，大多数人相信，如果他们的收入加倍，他们的幸福也会加倍。这是否符合事实？这项研究结果表明，事实上，人们的收入水平从每年2.5万美元提高到每年5.5万美元，收入提高1倍以上，而他们的幸福程度只增加了9\%，不到10\%。除此之外，关于这个主题有一项最广泛得到应用的研究，它告诉我们，如果你能得到一份稳定的中产阶级收入，在美国就是每年收入7.5万美元左右，之后挣再多的钱，对你个人的幸福水平也不会造成明显差别。

“那么，这又有什么意义？”你也许会问。

真相是，更多的研究证明，钱确实能让我们更加幸福。科学家研究表明，每天只要花上区区5美元，就能明显地提高你的幸福感。怎么会这样？关键不是你花多少钱，而是你决定如何来花这些钱。“每一天的消费选择释放出一系列的生活和情绪效应，这从唾液中能检测出来。”哈佛大学的伊丽莎白·邓恩和迈克尔·诺顿在2013年的著作《花钱带来的幸福感》中写道：“尽管有更多的金钱，能买来所有种类的奇妙物件，从更美味的食物到更安全的社区，但是金钱的真正力量并非来自数量大小，而是来自我们如何花这些钱。”

他们已经科学地证明了，你可以有很多不同的方式花你的钱，这样居然能显著提高你的幸福感。我不想在这里列举他们讲的那些更加幸福的花钱方式，希望你能好好地读他们的那本书。不过其中有三个最重要的花钱方式，我在这里一定要说一下：

第一，投资于体验，例如旅游、学种新的技能、学些课程而不是获取更多的财物，实不如虚。

第二，为你自己买时间。“如果我们把那些我们最讨厌、最畏惧的劳动外包出去，从刷马桶到清洁水槽，金钱能够改变我们花费时间的方式，把我们解脱出来，去追求我们的激情。”

第三，投资于别人。你猜对了。把你的钱给出去，竟然能够让我们真的感到很快乐！

研究表明，你给予别人的越多，你自己就越快乐。你拥有的越多，你就能给予越多。这其实是一个良性循环。邓恩和诺顿用他们的科学研究表明，把钱花在别人身上比花在他们自己身上更让人满足。这种好处不仅会延伸到主观上让你感觉更幸福，还会延伸到客观上让你的身体更健康。

换句话说，给予能让你更幸福而且更健康。

按照两位作者的研究，这种现象在全世界都很普遍，不管什么地域文化，不管国家是否富有，不管人群收入高低，不管年纪大小，“从一个加拿大的大学生给妈妈买一条围巾，到一个乌干达的女人给朋友买能够救命的治疗疟疾的药”。而且数据表明，礼物的价格并不重要。

他们的研究中有一项是，作者给实验参与者每个人5美元或者20美元，让他们一天之内花完。作者告诉其中一半人给自己买些东西，告诉另一半人用这些钱去帮助别人。“到了晚上，那些受到指派把钱花给别人的人，在情绪上感觉这一天过得更加幸福快乐，幸福感明显高于那些受到指派把钱花到自己身上的人。”他们这样写道。

这两个作者的同事、心理学家劳拉·阿克宁是西蒙弗雷泽大学的教授，她做了另外一项研究，向实验参与者每人发放价值10美元的星巴克星礼卡。

\begin{itemize}
\item 指导一些人单独进入星巴克，用星礼卡来给自己买东西。
\item 告诉一些人用星礼卡带别人去喝咖啡。
\item 告诉一些人把星礼卡送给别人，但是他们不准和那个人一起进入星巴克。
\item 告诉一些人带另外一个人和他们一起进入星巴克，但是只能用星礼卡给自己买咖啡喝，不能给和他们一起进来的那个人买咖啡喝。
\end{itemize}

这一天结束的时候，你认为哪些研究对象会说自己是最快乐的？如果你选择的是那个进入星巴克里面给别人买了一杯咖啡的人，你就选对了。按照作者的研究，能和自己帮助的人有沟通联系，“看到他们的慷慨行为产生的影响”，这样帮助别人的人是最快乐的。

我们从帮助别人这个行动上感受到的快乐，不但会更加强烈，而且会持续更久。我采访著名行为经济学研究专家丹·艾瑞里时，向他提出金钱和幸福这个话题，他告诉我：“如果你问别人‘什么会让你快乐，给自己买东西，还是给别人买东西’，他们会说‘当然是给自己买东西’。但是事实并非如此。研究表明，人们给自己买东西的时候，他们只能快乐几分钟，最多几个小时，但是如果他给别人买东西，即使买的是一份小小的礼物，这种给予的快乐至少会持续一整天，而且经常会持续好几天，甚至有时候持续好几个星期。

丹·艾瑞里也告诉我一个“美好心灵实验”，有一家公司的雇员每个人都得到3000美元的奖金，有的人把奖金花在自己身上，有的人接受劝导把钱捐出去了。猜一猜，谁更快乐？

“6个月之后，那些把钱捐出去的人说自己快乐的程度大大超过那些把钱留给自己的人。”丹·艾瑞里说，“我的意思是，我们应该好好地想想给予的意义是什么，对不对？给予这个事情非常了不起，能够把你和世界的其他人联系在一起，这是一个循环，你给予别人好处，你从中得到了好处。

你把钱给出去了，特别是你用钱为一个陌生人做了一件好事，相比于你给你所爱的人做了一件好事，这种幸福放大的程度是呈几何级的，成倍甚至数倍放大。花钱给你爱的人做好事相当于你现在的薪水，花钱给陌生人做好事，相当于你的薪水加了两三倍。

基于我个人的经验，我亲眼看到，当你给予的时候，很多令人惊奇的事情就会发生在你身上。当你超越了自己的生存，超越了你只顾追求自己成功的心态，达到一个新的精神境界，你活着远远不止为了自己，突然间，你的恐惧、你的懊恼、你的痛苦、你的不快乐会通通消失。我真的相信，当我们奉献出自己的时候，生命、神明、荣恩（不管你怎么称呼它，意思都一样）就开始介入，指引我们。记住，生命支持那些能支撑更多生命的生命。

我给你举一个例子，一个小男孩的人生怎么被重新点燃。在此之前，他在康涅狄格州牛顿镇，经受了恐怖的校园袭击，他的心灵和灵魂几乎完全被摧毁了，这个故事讲述的是，通过给予的行动，找到人生目标，得到灵魂上的启发，从痛苦中解脱。

\subsection{超越痛苦的力量}

杰特·刘易斯永远不会忘记2012年12月14日。那天早上，一个精神失常的人举着枪闯进桑迪胡克小学，开枪杀死了26个人，其中包括20个5\textasciitilde{}10岁的孩子，然后开枪自杀。就在他狂暴开枪射击的过程中，有一个时刻，杰特6岁的弟弟杰西注意到枪击者的枪卡住了，于是对班里同学大喊快跑，结果很多同学跑了出去，这个勇敢的小男孩那天早上挽救了很多人的生命，但不幸的是，他没有挽救自己的生命。枪击者掉转枪口，对着杰西，连开几枪，将其击毙。

想象一下，如果杰西是你的儿子或者你的兄弟，你会心碎成什么样子，你会受到多大的伤害。我有幸见到了13岁的杰特，还有他和杰西的妈妈斯嘉丽，那是在我飞到牛顿镇参加这个重大枪击事件一周年纪念仪式的时候，我想帮助这群幸存者应对这场毁灭性悲剧带给他们人生的持续影响。正如我所预料的那样，这些家庭中的很多人深受痛苦悲伤的折磨。但是我和杰特交谈的时候，我非常吃惊地了解到，他的痛苦和折磨只是在和一群非同寻常的卢旺达孤儿交流了一次之后，就彻底改变了。这些年轻的卢旺达男孩和女孩听说了杰特经受到的惨痛经历，伸出援手，希望能够和杰特网络视频通话，跨越半个地球，分享他们治愈巨大痛苦的感悟。

这些卢旺达孤儿都经受了历史上最大的悲剧之一。1994年，卢旺达发生种族大屠杀，导致多达100万名图西族人死亡，杀死这些人的凶手是他们的邻居胡图人，这场大屠杀持续了100天左右。这些幸存下来的孩子主动联系杰特，用网络电话视频交流。其中有一个叫香塔尔的女孩告诉杰特说，她为杰特失去他的弟弟感到非常难过，但是她想让杰特知道，能夺走你人生幸福和快乐的只有你自己，那个枪击者并没有那种能力。

香塔尔接着分享了她自己的亲身经历：她只有8岁的时候就经历了极其恐怖的一幕，一个胡图人用大砍刀砍死她的父母，接着那个凶手又转向她，用刀猛砍她的脖子，把她瘦小的身体，扔到一个巨大的集体埋葬墓地。身体埋到土里，伤口鲜血直流，惊恐万分，但是她心里充满了要活下去的愿望。香塔尔用她的双手，从那个埋得较浅的坟墓中刨出一个洞，爬了出来，跑到她们村庄上面的大山里，得以重获自由。她躲在黑幽幽的森林里，往下能看到她原来称之为家园的村庄，大火吞没了一座又一座房屋，空中回响着她熟悉又热爱的人的尖叫声。香塔尔靠吃草活了一个月，一直等到那场大屠杀停止下来。

你会想，一个孩子被迫见证了父母被人砍死，肯定会在情绪上受到惊吓，持续一生不能解脱。你会预料，这个8岁的女孩从此一生都会生活在愤怒和恐惧中。但是她没有，因为她掌握了能够塑造我们人生的三项重要决策。

香塔尔告诉杰特：“我知道你现在不会相信，但是你马上就能治愈，过上快乐和美好的人生。只要你训练自己，每天都训练自己，要感恩，要宽恕，要有同情心。对你所拥有的东西要感恩，不要关注那些你并未拥有的东西。你必须宽恕枪击者和他的家庭，找到一种方式服务他人，这样你才能从痛苦中解脱出来，心灵重获自由。”杰特从电脑视频通话的屏幕上看到，香塔尔的脸庞充满着一种快乐，那种快乐远远超过杰特的想象，尽管杰特的人生同样非常悲惨，但是这个女孩描述的她经历的那些恐怖情景，远远比他能够想象的恐怖。这个女孩都能够从自己巨大的痛苦中解脱出来，我应该也能。现在是时候开始行动，寻求解脱了。

但是，他该怎么做？杰特心想，他必须找到一种方式来回报这个年轻的高尚的灵魂，是她从千里之外，主动联系他进行网络视频通话，在自己需要的时候，送来温暖和爱的。香塔尔找到了活下去的理由，找到了激情。找到了人生目标的她决定去保护，去爱，去扶养这场种族大屠杀中一些比她更小的孤儿。这成了香塔尔的人生使命，让她不再只关注自己，只关注自己失去亲人的痛苦，从而心灵从升华中得到解脱。

香塔尔这种服务他人的事迹深深触动了杰特，他开始执迷于给予他人、帮助他人、奉献他人的想法。杰特决定要帮助这位非凡的女孩创造更好的未来，这就是他的使命。杰特开始日夜工作筹集资金，想让香塔尔有足够的钱读完大学。过了几个月，这个13岁的男孩学会用网络视频电话联系香塔尔，骄傲地宣布他已经筹集了2100美元，足够送香塔尔去读一年的大学！香塔尔简直不敢相信，她太感动了。但是像很多年轻人一样，尤其是在第三世界国家，上大学对于她来说并不是一个很有实效的选择，特别是她已经开始创办自己的小企业——一家小小的商店。（你也许完全可以想象，有如此精神的女人，肯定会是一个非常成功的企业家！）于是为了继续发扬光大这种给予的精神，香塔尔把这份令人震惊的礼物转送给她最好的朋友贝蒂，贝蒂也是一个孤儿，上一次和香塔尔一起打电话鼓励杰特。

杰特这种爱的无私奉献，深深感动了我，我当场决定为贝蒂提供其余三年的学费，同时也大力支持香塔尔，为她提供资金建造一个全新的商店，也建造一个永久性居所，让她扶养的那些孤儿有足够的地方居住。

今天，我们正在一起努力工作，扩大、募集更多的资源，有7.5万个孩子在那场种族大屠杀中幸存下来，我们希望能够帮助更多这样的孤儿。\footnote{\label{fn:pdf-46-1}我们训练心理学家和职业教练，让他们学会实际操作技巧的和心理技巧，在这些危机发生的时候，能够参与其中，发挥作用。如果你的条件符合要求，又愿意在危机期间参加志愿工作，请联系安东尼·罗宾斯基金会（www.anthonyrobbinsfoundation.org）。}从杰特和香塔尔两个孩子身上，我们可以学到的经验是：人类可以克服我们自身的痛苦，只要我们选择看到生命的美好，找到办法给予和奉献自我。你给予别人帮助，上天就会给予你治愈痛苦创伤的礼物。关键是找到某个东西，能激发你愿意给予、愿意付出。这种使命感，正是生命的终极力量。就在这个时候，你真正变得富有了，因为你从一种只是贪图享受的人生，转移到既快乐又有意义的人生。

\subsection{给予就是治愈}

当然，给予并不仅是给予金钱，也可以是给予你的时间、给予你的情绪、给予你的存在，陪伴你的孩子、陪伴你的家人、陪伴你的爱人、陪伴你的朋友、陪伴你的同事。我们的工作也是我们给予别人的礼物，不管这个礼物是一首歌、一首诗、一个跨国企业，还是医护人员、教师、咨询顾问，我们都有某些东西可以给予。事实上，仅次于爱，我们可以给予的最神圣的礼物就是我们的劳动。当个志愿者，义务劳动付出你的时间，给予别人只有你才能做到的细心关照，分享你的技能，也会给你很大的“回报”。

我的朋友阿里安娜·赫芬顿写了一本精彩的书《发展》，她在书中引用的研究展示给我们，给予的行动如何能实实在在地改善你的身体健康和心理健康。其中我特别喜欢的案例是，英国大学埃克塞特医学院2013年做的一项研究表明，志愿工作通常会降低人的抑郁水平，增加幸福感，使死亡率下降22\%！她还写道：“志愿工作之后，至少会有一个星期的时间让你的幸福感提高，相当于你每年的工资从2万美元提高到7.5万美元带来的幸福感。”

那么，财富的最终秘密是什么？这就是任何形式的给予，都能给你增添更多的财富，远远高于得到为你增加财富的速度。我并不关注，我们作为个人而言是多么强大有力，不管你是一个企业大亨、政治领导人、金融巨头，还是娱乐巨星，过上一个圆满人生的秘密，并不只是把事做好，而且要多做好事。毕竟我们都知道这样的故事——那些超级富豪，有一天突然醒来，意识到：“生命并不只是我一个人的事。”

\begin{BookQuote}
成为躺在坟墓里的最富有的人，这件事对于我来说并不重要。晚上上床能说我们今天做了一些很棒的事，这对于我来说才是很重要的。

\hfill ——史蒂夫·乔布斯
\end{BookQuote}

在19世纪之前，大多数慈善活动都是由宗教组织举办的，直到后来钢铁巨头安德鲁·卡内基开始大做慈善活动，慈善事业才发生巨变。原来那些国王和贵族，还有最富有的家族，并没有兴趣回馈社会：对于大多数有钱人来说，他们只想紧紧攥住自己手里的钱，只留给自己，或者死后留给自己的继承人。很多企业界人士也有同样的想法。但是卡内基领导和他同时代其他“强盗企业大亨”，创造出了我们现在这样的慈善事业。

卡内基是一个严酷无情的企业家，但他制造出来的钢铁，用来修建铁路和摩天大楼，改变了整个美国。卡内基先给社会增加价值，自己从中赚到了钱，这样整个社会得到了好处，他也从中得到了好处。后来卡内基成了他所在时代中最富有的人。卡内基随之就来到这样一个人生阶段：所有他想要得到的东西，他都得到了，而且比他想要得到的还要多得多。卡内基拥有的钱实在太多了，他开始认识到，钱本身并没有多大意义，除非他能够用这些钱来做一些超越他自己的事情。于是卡内基用前半生给自己赚钱，用后半生给别人花钱。卡内基写了一篇文章，描述自己的人生观、财富观的彻底转变，现在还非常值得阅读，这篇文章的题目是《财富的福音》，这篇文章后来被扩展成一本书。我的朋友、诺贝尔经济学奖获得者、耶鲁大学经济学教授罗伯特·席勒坚持要求他所有的学生都必须好好读这本书，因为他想让学生知道，资本主义也可以成为一股善的力量。卡内基的这篇文章改变了社会，影响了与他财富相当的富豪，甚至挑战了他最大的竞争对手，财富巨大到不可思议地步的约翰·洛克菲勒。受到这个凶猛竞争对手精神力量的感召，洛克菲勒开始把大量的金钱捐给美国一些优秀的慈善基金会。卡内基创造了一个新的人生衡量标准：这个标准衡量你个人的重要性，不是根据你拥有了什么，而是根据你给予了什么。卡内基的捐赠，让美国的图书馆数量增加了1倍，为社会发展提供了如此之多的智力资本，促进了我们的社会发展，只有后来出现的互联网才能媲美卡内基的巨大贡献。

我们的朋友查克·费尼成了“现代卡内基”，他把自己拥有的75亿美元的个人财富几乎全部捐出，只不过他选择沉默，绝口不提自己的捐赠，他的事直到最近才被披露出来。

我第一次见到查克·费尼的时候，他已经83岁了，进入了人生的最后阶段。他想说话时间长一点儿都很困难，但是他一出现在你面前，你就会产生一种言语无法描述的体验。他在你面前，你就能感觉到一种人生活得非常圆满的力量。你可以看到，从他眼中散发出来的快乐，从他脸上随时闪现的微笑，从他内心散发出来的友善。

查克·费尼在卡内基之后又影响了一代人。很多人说，特德·特纳是下一个重新点燃这种大规模慈善捐赠的人，他承诺要捐赠10亿美元给联合国。从那以后，比尔·盖茨和沃伦·巴菲特联合发起了给予誓言活动，激励世界上的富豪至少把一半的个人财富捐赠给慈善事业。按照最近的统计，超过120个亿万富翁已经签署了给予誓言，包括本书中我采访过的一些超级富豪，例如瑞·达利欧、布恩·皮肯斯、萨拉·布雷克里、卡尔·伊坎、保罗·都铎·琼斯。

布恩·皮肯斯告诉我，他的慈善捐赠有点儿过头了。他最近捐赠近5亿美元给母校俄克拉何马州立大学，这让他的捐赠总额超过了10亿美元。可是他最近的投资出现了一些亏损，让他的个人财富下降到只有9.51亿美元，自己手里有的钱竟然没有捐的多，这成了一大奇闻。但是布恩·皮肯斯并不介意。毕竟，他才86岁。“不要担心，托尼。”他说，“我计划未来几年再赚两个10亿美元。”他根本没有损失吃亏的感觉，他从自己的给予行为中得到的快乐是无价的。

在现代社会，全世界最富有和最有影响力的人解决了世界上最大的问题。卡内基捐赠图书馆来解决教育问题，比尔·盖茨和梅琳达夫妇捐赠大量资金用于设立奖学金，并发放疫苗预防传染病。波诺的激情是豁免第三世界国家的债务。但是你一定先要成为一个亿万富翁或者一个摇滚歌星，然后才能解决世界上最大的问题吗？在现在这个互相联结的世界中，并不需要像这样先有钱或者先有名才能解决世界上的大难题。我们可以运用科技的力量一起工作，一起努力，我们每个人都做出自己的一点奉献，但是它们合在一起仍然能够产生巨大的影响力。

\subsection{消除饥饿，消除疾病，消除奴役}

我不能确定你的激情是什么，但是有个领域——我个人曾经有过类似的经历，能够深切地感受到那些人受的苦和受的罪，就是那些急需帮助的孩子和家庭。你甚至需要把冰块放在你的血管上，以免感受到那些孩子的痛苦让你过于激动。所以，我们花上一分钟来看看现在折磨孩子和他们家庭的三个最大的问题，我们可以采取三个能马上见效的具体步骤来发挥我们的影响力，改变这些家庭的现状。

第一个问题是饥饿。你认为，在美国这个世界上最富有的国家，还有谁会每天晚上上床睡觉的时候饿着肚子？根据美国人口统计局的调查，结果听起来令人震惊，事实上，美国5岁以下的孩子里，每4个孩子里面就有1个生活在贫困中，10个孩子里面就有接近1个生活在极端贫困之中。（贫困的定义是年收入低于11746美元，或者每天收入低于32美元，一个四口之家只能靠这些钱来维持生活。）

美国有5000万人，其中包括近1700万个孩子，生活在食物得不到保障的家庭里，或者按照纽约对抗饥饿联合会的乔尔·博格告诉莫耶斯公司的特丽萨·莱利的那样，有些孩子所在的家庭“没有足够的钱来定期购买他们需要的食物”，有些孩子生活的家庭“只能定量分配食物，吃了上顿没下顿，父母没有食物来喂养自己的孩子”。与此同时，美国国会决定每年削减87亿美元的补充营养协助计划福利，以前称之为食物券，相当于让50万个美国家庭每个月少了一个星期的饭钱。

我小时候就生活在这样食物短缺的家庭中。也正是因为有过这样的经历，我特别想在这个领域发挥影响力，做些实实在在的事情，给那些孩子的生活带来一些改变。我知道这不仅是统计数据，那些数字背后都是活生生的人天天饿着肚子，他们非常痛苦。

我曾讲过，我的人生之所以发生改变，是因为我11岁时有个陌生人给我们家送来了一篮子食物。这件事能改变我的人生，并不只是因为我们收到了食物，还因为有个陌生人关心我们。这个简单的行动产生了指数型的影响，我一直在回报我得到的那份感恩节礼物，38年间，我给4200万人提供了食物。关键是我没有等到我有能力大规模处理这个巨大的饥饿问题，才开始行动。我并没有等到我变得很富有，才开始行动。我马上就着手解决这个饥饿问题，尽管当时我手里只有一点点钱，但是我能做多少是多少，关键是我马上行动起来了。

一开始，我的财力非常有限，只够给两个家庭提供食物，但是后来我受到了激励，我把目标提高了2倍——给4个家庭提供食物。第二年，我的目标提高到8个家庭，再过一年又提高到16个家庭，随着我的公司规模不断扩大，我的影响力也在增长，我捐助食物的目标也大幅度地提高，变成捐助食物给100万个家庭，后来提高到200万个家庭。就像投资收益按复利增长一样，我在给予上做的投资，得到的回报也是按复利增长，而且增长得更多。我很荣幸现在能够捐助5000万份餐食，有了你和其他人一起合作，能够提供超过1亿份餐食，能有这种荣幸，我内心的幸福感简直无法形容。过去我是一个接受别人捐赠食物的孩子，承蒙上天的恩赐，加上个人的努力，现在我很荣幸能够给别人提供食物，当初别人行善捐赠一篮子食物给我和家人，现在我能加倍行善给上亿人。

没有什么能比得上人类的灵魂激情燃烧的力量。一路上，别人的关怀触动了我，也有许多好书触动了我。是这些好书，把我从一个受到限制的世界，带入一个充满无限可能的人生，因为我在读书时仿佛进入了作者的大脑，看到他们如何改变了自己的人生。遵循这种传统，我去找了我的出版商西蒙–舒斯特公司的编辑，让他们知道，我不仅给身体提供食物，还给思想提供养分。

但是现在我想请你考虑一下，和我一起合作，用一种方式，可以在接下来的几年里继续做这些善事。这是一个很简单的策略，能够提供100万份餐食，不仅是今年，而是每一年，给那些需要食物的家庭。这并不需要你做出大额的捐赠。我提议的这个捐赠计划，给你提供能够改变生命和挽救生命的机会，只需要毫不费力地捐赠一些零钱即可。如何操作？跟我一起加入三大战役：消除饥饿，消除疾病，消除奴役。

\subsection{用你的零钱来改变世界}

我给你一个提议。我写本书的目标是，帮助你区别洞察力和技能，给你一个计划，能够真正赋予你能力，为你和你的家庭创造出长期持续的财务安全、财务独立、财务自由。我执迷于找到给你的人生增加更多价值的方法，比你想象的还要多得多，只用一本书就能做到（尽管我必须承认这是一本很厚的书）。我希望本书能激发你超越短缺，实现成为一个富有的人的意愿。你成为富人的那一天，就是你开始内心充满快乐地给予的那一天，不管你在财务上有多少财富，你在精神上已经非常富有了，但是给予必须是发自内心的给予，并不是出于你必须去做才给予，也不是出于愧疚感才给予，更不是出于别人的要求去给予，而是出于给予能让你内心兴奋激动。

根据美国劳工部的统计数据，美国1.24亿个家庭中每年在娱乐上的平均花费有2604美元，那么美国人只是在娱乐上的花费总额就超过3200亿美元，想象一下，要是能从这3200亿美元里拿出一部分解决以前非常棘手的大问题，例如饥饿、贩卖人口、缺少清洁饮用水，结果会怎么样？在美国，只需要1美元，就能给那些饥饿的人提供10份餐食。想象一下，要每年提供1亿份餐食，那会帮助多少人！但这只需要花0.1亿美元，只相当于我们每年花费在娱乐上的3200亿美元的0.0034\%，就是1\%的1/3！这相当于1美元里拿出来3美分多一点儿！于是我跟一些商业界和市场营销界的杰出人士合作，其中包括社会资本家鲍勃·卡鲁索（全球排名前100的对冲基金的高桥资本管理公司的前任管理合伙人和首席运营官），还有我的好友、慈善家马克·贝尼奥夫（Salesforce.com网络公司创始人兼首席执行官），我们共同开发出一种技术，能让你轻松地、毫无痛苦地拿出这些零钱来挽救生命。

只用不到一分钟，你就可以上网选择参加消除计划，这样每次你在世界各地使用你的信用卡时，你的购物消费支付价格会自动取整到最接近的美元数字，那么这个不足1美元的零头就是你的小小捐赠，它会直接转到批准的有效慈善活动中，以后我们会向你汇报你的捐赠让那些受助者发生了巨大变化的故事。我来举个例子说明这是怎么运作的。例如，你在星巴克买了一杯咖啡，价格是3.75美元，那么如果你参加了这个捐赠计划，你刷卡支付的价格会自动取整为4美元，相当于你多付了0.25美元的零钱，它会按照设定好的路径捐赠给预先选择好的慈善机构。对于每个消费者来说，所有这些化零为整捐赠的零钱，一个月下来平均每个人不会超过20美元。你可以给你的这种化零为整的零钱捐赠设定一个每月捐赠总额的上限，但是我们建议你每个月捐赠的上限不要低于10美元。

你想知道你这样化零为整的零钱捐赠，会有多大的影响力吗？比如，你一个月捐赠零钱达到20美元：

\begin{itemize}
\item 你每月捐赠零钱达到20美元，可以为饥饿的美国人提供200份餐食（也就是说，一年下来你一共能够提供2400份餐食）。
\item 你每月捐赠零钱达到20美元，你可以提供一份清洁的、可持续的水源，能满足10个印度孩子每月的饮水需要，这样的话，每年你一个人就能够保护120个印度孩子避免患上饮用水污染所带来的疾病。
\item 你每月捐赠零钱达到20美元，就可以把这笔钱作为一笔首付款，用它来解救一个因走私交易而变成奴隶的柬埔寨年轻女孩，让她重获自由。
\end{itemize}

这是全世界孩子和家庭面临的三大难题。在美国，最主要的问题是饥饿。也正是因为这个原因，我们专注于消除饥饿，因此我们和消除美国饥饿组织一起合作，共同努力。

但是从全世界来看，孩子们面临的最大挑战是疾病。你知道吗？由于水污染引发的疾病，是人类的头号杀手，每年导致340万人死亡，这是世界卫生组织的官方统计数据。事实上，每隔20秒，就会有一个孩子死于水污染导致的疾病，水污染导致的疾病“杀”死的人数，比“二战”以来所有因武装冲突死亡的人数还要多。

这也正是消除计划的第二个承诺：消除水污染传播的疾病，提供清洁的饮用水给全世界尽可能多的孩子。有不同类型的组织机构，提供不同类型的可持续性解决方案。有些解决方案只要求每人2美元，就能提供给这些孩子和他们的家庭稳定可靠的清洁饮用水。

\subsection{自由的代价是多少？}

综观全书，我们一直在努力，以确保你能获得财务上的自由。那么只投资相当于你每月支出的一小部分，来帮助一个孩子重获人身自由，你觉得值不值？每年都有840万个孩子被拐卖成为奴隶。2008年，美国广播公司的新闻记者丹·哈里斯暗访，描述了需要多长时间和多少钱就可以买一个小孩当奴隶。这个记者离开纽约，经过10个小时到达海地，讨价还价一番，只花150美元就买到一个小孩当奴隶。这个记者说，在这个现代世界，买个小孩当奴隶，比买一个苹果音乐播放器还要便宜。

你肯定无法想象，这种事要是发生在我们自己的孩子身上，或者发生在我们喜爱的任何一个孩子身上会怎样。但是你要努力想象一下，你的捐赠行动，让一个孩子的人身重获自由，让受到多年奴役的灵魂重获自由，会是什么样的。你的善行，能给那个当奴隶的孩子带来多大恩赐，根本没有语言能够形容。而且，你知道，即使在你睡觉的时候，你的捐赠也在帮助那些每天正在努力赢得这场战争的人。

那么，我们应该如何处理这些巨大的挑战？我们每个人都要一起努力，一次解决一点点，一次改善一点点。今年你和我，还有我的一些朋友，会努力给1亿人提供餐食。但是要是能用一种可持续的方式，每年都给1亿人提供餐食，是不是有点儿太不可思议了？我每天给10万个印度人提供清洁的饮用水，这是我很有激情、非常热爱的慈善事业之一。要是我们一起努力每天能给300万人提供清洁的饮用水，而且捐助人数继续增长，是不是一件很令人震惊的事？或者我们一起努力解救出5000个受到奴役的孩子，支持他们接受教育，让他们找到途径重新开始健康的生活，怎么样？

这就是我们只要有10万人参加消除计划，合在一起的力量就能够做到的事。就像我建立了一个个人基金会以后，这个使命就呈几何级增长，如果再过10年或者更长一段时间，我们就可以找到一种方法，把参加消除计划的人数增加到100万人，那么我们就能够每年提供10亿份餐食，给3000万人提供清洁饮用水，把5万个孩子从被奴役中解救出来。这些数字都大得惊人，但是生命无价，即使只能够挽救一个孩子的生命，那也值得我们做所有这些努力。

那么你的愿景是什么？大多数人过高地估计他们一年之内能做成什么事，而经常过低地估计他们10年或者20年能做成什么事。

我可以告诉你，当我开始执行我的个人使命，给两个饥饿家庭提供食物时，我很激动。后来我的目标变成给100个家庭提供食物，后来这个目标提高到1000个家庭，再后来提高到10万个家庭，再后来提高到100万个家庭。我们自己做得越大，我们看到自己有可能做到的潜力也越大。这取决于我们自己。和我一起努力吧，拿出来你的一点儿零钱，我们一起改变这个世界，消除饥饿，消除疾病，消除奴役，让世界上更多的孩子生活得更美好。

\begin{BookQuote}
我发现，给予能带来很多好处，最大的好处是解放了给予者的灵魂。

\hfill ——玛雅·安吉罗
\end{BookQuote}

不管你有没有参加消除计划或者其他慈善活动，做出一个决定，从你赚的钱里面拿出来小小的一部分，或者拿出来一小部分时间，有意识地选择投资到某件事上，这件事并不会让你个人直接受益，而是帮助某些急需帮助的陌生人。这个决定不涉及对错，而关系到你真正的财富，让你真正地感到更有活力，内心真正地感到满足。

在《花钱带来的幸福感》这本书里，邓恩和诺顿写道：“给予别人，给予和你无关的陌生人，在做对了的时候，你感觉像是一个选择，把我们和他人联系在一起，此时就产生了一种影响，即使是小礼物也能够增加幸福感，这有可能引发多米诺效应，激起一连串的慷慨捐赠行为。”

做这种“有利于社会的消费支出”（比如送别人礼物，捐赠做慈善活动），影响力很大，这触动了丹·艾瑞里和他妻子两个人，他们决定付诸实践，制定一个简单的花钱制度，让他们夫妇两人和两个孩子作为一个家庭能够一起坚持把慈善活动做下去。孩子们得到零花钱，必须把钱分成三份，分别放在三个罐子里。

第一个罐子里的钱，花给自己。

第二个罐子里的钱，花给他们认识的某人。

第三个罐子里的钱，花给他们不认识的人。

请注意，这三个罐子有两个罐子里的钱都是花在别人身上的，都是有利于社会支出的，因为这样花钱能让孩子感到快乐。把所有钱分到三个罐子，分别花到自己、熟人、陌生人身上，这些钱都是应该花的，但是丹·艾瑞里非常细心，把钱分成相等的三份，花到他们所不认识的人身上的那一份，和花到自己身上以及花到朋友和家人身上的钱是相同的。花钱送给朋友和家人礼物当然会感觉很好，因为这是给予你所爱的人的礼物，但是用第三个罐子里的钱来做慈善，把钱花到陌生人身上，可能是最让人感到满足的，也是最重要的给予方式。

有些人尽管自己的日子过得有些困难，再给予别人很不容易，但是还是想方设法地给予别人，他会有非常大的善报。给予别人，多做善事，会对我们的大脑产生启动效应。这样做会训练我们，让我们逐步养成习惯，知道并不是做到现在这样就够了，我们还能做得更好。只要我们的大脑相信了，我们就能做到，我们就能体验到。

约翰·邓普顿爵士并不只是世界上最伟大的投资人，也是世界上最伟大的人，在大约30年前，他跟我分享了一些事：他说在他这辈子认识的人里面，有些人能坚持从自己的收入里拿出8\%\textasciitilde{}10\%捐赠给宗教机构或者慈善机构，凡是这样每年捐赠1/10的收入而且连续坚持10年的人，据他所知，没有一个人不是财富大幅度增长的。但是问题就在这里，每个人都会说“等我日子过得好些了，我会给予的”。我过去有段时间也是这样想的。但是现在有一点我可以做证：不管你现在日子过得怎么样，都应该马上开始给予。你应该马上开始培养这种给予的习惯，即使你觉得现在还没有准备好，即使你觉得根本没有多少剩余的钱去给予。为什么？因为正如我在本书第1章里跟你说过的那样，如果你有1美元的时候没有给予别人10美分，那么你有1000万美元的时候，你就不会给予别人100万美元。

你如何扩大你给予别人的遗产？你如何给予别人你的时间和能量？你愿意拿出你个人收入的1/10捐赠给慈善事业吗？你愿意拿出一分钟上网签约消除计划，把你的零钱变成投资改变那些孩子的人生吗？如果你受到激励，请你现在就这样做，你就联结上了你能拥有的影响力。记住：从你的给予中受益最多的人很有可能是你自己。一个慈善家的人生，也是从简单的一小步开始的。让我们一起迈出这小小的一步吧。

\begin{BookQuote}
我不去想所有的悲惨，而去想还剩下来的那些美好。

\hfill ——安妮·弗兰克
\end{BookQuote}

顺便说一下，我并非一开始就充分意识到感恩和给予的意义。我过去也是一直过得紧巴巴的，缺这又缺那。回头看，我的人生并不是一直过得一帆风顺，但是我的人生总是得到上天的赐福，可是那时候我并没有意识到。因为在我长大的过程中，从钱上来讲一直是很贫穷的，我一直努力工作，以确认我能够获得最高水平的成就，但是我并没有认识到成就是一阵阵爆发而来的，不可能时时都处于最高水平。

要花很长时间才能真正学会一样东西，而且还要花很长时间才能达到精通，到了这个水平，这样东西在你脑子里就变得根深蒂固，成了你人生的一部分。所以我刚刚开始时，遭受到一系列的挫折。我如何应对这些挫折？我得实话实说，当时我可没有那种开悟之后的智者为人处世的优雅！我一直很生气，很懊恼，简直气坏了！没有一件事顺心如意。我的钱快要花光了！

有一天午夜，我开着车行驶在57号高速公路上，接近坦普尔大街出口匝道，快到加利福尼亚州的波莫纳市了，我心里一直在嘀咕：“我什么地方做错了？我工作这么努力，我到底缺少什么？为什么我失败得这么惨，无法得到我想要得到的东西？为什么我这么努力却不行？”突然之间，我两行热泪直流，我把车停到路边，翻出来我每天随身携带的日记本，借着仪表盘的灯光，开始狂乱地书写。我用大字在整整一页上写下这句话给我自己看：“生活的秘密是给予。”

是的！我意识到我忘了生活究竟是为了什么。我忘了正是在这里才能找到所有的幸福快乐，因为生活不仅是为了我自己，生活是为了我们。

我开车回到高速公路上，我受到了激励，重新开始聚焦，重新意识到我的使命，这让我再次激情燃烧。我开始有一阵子做得挺好，但是不幸的是，那天晚上我所写的东西只是一个概念而已，其实这只是表面上的领悟，还完全没有体现在我的具体行动上。接下来，我开始碰到更多的困难和挑战，6个月之后，我的钱全部亏光了。不久之后，我发现自己正处在人生的最低谷，一个人住在一个只有37平方米的单身公寓里面，就在加利福尼亚州的威尼斯市，我内心充满怨恨。其实只要你追求的是相当大的目标，你就会遇到这些困难和挑战，这是自然而然的事，我却为此指责所有的人，没办法，当时我陷入这种误区出不来了。我断定我被那些人利用了，他们就是想占我的便宜。“如果不是因为他们，我肯定现在发展得很好！”我对自己说。于是我到处抱怨求安慰。我变得越来越愤怒，越来越懊恼，于是我就变得越来越低产出、低效率。

心情郁闷，我就开始胡吃海喝，这是我逃避现实的方式，吃的喝的都是那些很糟糕、很可笑、根本没有营养价值的快餐。才过了几个月，我的体重就增加了37斤多。这可不容易，你得吃好几千斤的食物，而且不能运动太多，才能这么快就长出这么多肥肉！我发现，自己现在做的事都是以前我嘲笑别人做的事——我不是在吃东西，就是在看电视连续剧。

回过头去看一看，当时的我堕落得是多么厉害，很可笑，当然也有点儿丢人。最后我只剩下19美元，还有几个硬币，我一点儿希望也没有了。当时我对一个朋友特别生气，我混得好的时候，他找我借了1200美元，却一直没有还钱给我。现在我手头很紧去向他要这笔钱，结果他根本不理我。我打电话给他，他竟然不接！我很生气，我心里慌张地想：“我的天哪，我该怎么办？我连吃饭的钱都没有了！”

不过我这个人一直是很务实的。我想：“好了，当初我17岁，无家可归，那时的我比现在惨多了，我是怎么熬过来的？”那时我去找了一家专门做自助餐的餐厅，你想吃多少就可以吃多少，只要花一点点钱就可以了。这让我有了一个好主意。

我住的公寓不远处有个美丽的海湾，叫玛丽安德尔湾，洛杉矶那些有钱人都把他们的游艇停靠在那里。那里有家餐厅叫作小公牛餐厅，自助餐非常丰盛，一顿只要6美元。我可不想把任何一分钱浪费在汽油上或者停车费上，于是我不开车，走了5公里路来到餐厅。这个餐厅正对着海湾，风景很美，我找了靠窗的位置坐下来，装了一盘又一盘的食物，猛吃猛喝，就像我活不到第二天了一样，也许可能真的会这样。

我吃饭的时候，看着那些游艇在海湾里来来往往，梦想着那种有钱人开游艇的日子会是什么样子。我的心态开始变了，我可以感到一层一层的愤怒被化解掉了。我快吃完的时候，忽然注意到餐厅外面有个小男孩走过来，穿着一套小西服，看上去只有七八岁，紧跑几步上前推开门，让他年轻的妈妈先进去。然后非常骄傲地领着妈妈来到他们的桌子前，帮妈妈把椅子拉出来，请妈妈坐好。他表现得特别有绅士风度。这个小家伙看起来非常纯真，非常善良。他是一个完完全全的给予者，他对待自己的妈妈的方式那么尊敬、那么热爱，你一下子就能看出来。他深深地打动了我。

我买了单，走到他们的桌子旁，我对那个男孩说：“打扰一下，我觉得你是一位非常优雅的绅士，我很想认识你。看到你能这么绅士地对待这位女士，我非常震惊。”

“她是我的妈妈，我当然要好好对待她了。”小男孩很自然地说。

“哦，我的天哪。”我说，“这就更酷了，你能请你妈妈吃午餐，你太厉害了！”

小男孩定了定睛，用平静的语气说：“哦，其实不是我请妈妈吃饭，因为我只有8岁，我还没有工作呢，无法赚钱请妈妈吃饭。”

“不，你可以请你妈妈吃午餐。”我说。就在那一刻，我把手伸进钱包，拿出我所有剩下来的钱，大概加在一起有13美元，还有一些硬币，我把这些钱全都放到桌子上。

小男孩抬起头来看着我，他说：“这钱我不能要。”

“你当然能要了。”我对小男孩说。

“为什么？”

我看着他的脸，给他一个大大的微笑，我说：“因为我比你年纪大呀。”

他两眼直盯着我，惊呆了，然后开始咯咯地笑。我只是转过身去，走出了餐厅。

我其实不是走出餐厅那扇门，我简直是一路飞奔回家里！我肯定是完全疯掉了，因为当时我名下一分钱都没有了，但是相反我感到完全自由、完全解脱了。

就是在那一天，我的人生永远改变了。

就是在那一刻，我成了一个富有的人。

我内心有某种东西最终超越了短缺的感觉，我从那些大家称为钱的东西里面解脱出来。我能把每一件东西都给予别人，内心没有一丝恐惧。有某种东西超越了我的思考，在我的精神深处，我知道我得到了上天的指引。碰到这个时刻，我们大家都知道。这个时刻是命中注定的，就像你命中注定现在会读到本书的这些文字。

我认识到，我一直都忙着努力得到，但是我忘了给予。不过现在我重新找到了自我，我重新找到了我的灵魂。

我扔掉了我的那些借口，不再指责别人，突然我一点儿也不生气了，我也不感到懊恼了。你也许会说，我简直太傻了！因为我根本不知道明天的早餐在哪里，明天的晚餐在哪里。但是我心里根本想都不会想这种事。相反，我内心感到无比快乐，让我从噩梦中解脱出来——别人做了那么多让我觉得我的人生注定失败的噩梦。

那天晚上，我下定决心实施一个重大行动计划。我设计得清清楚楚，我要去干什么，如何让自己重新找到工作。我觉得我肯定能实现计划，但是我并不知道我什么时候能拿到第一笔工资，甚至比这更加紧急的是，我下一顿饭钱在哪里？

奇迹发生了。第二天早上，有一封老式的传统信件送到我楼下的信箱里。那是我的一个朋友手写的信件，他在信里说他非常抱歉，一直回避不接我的电话。过去他需要我的时候，我都会出手相助。他知道，我现在肯定是遇到麻烦了。所以他把向我借的那1200美元都还给我，还加上了一点儿钱作为利息。

信封里有一张支票，价值1300美元。这么多钱，足够我花上一个多月了！我哭了，一下子如释重负。然后我想：“这意味着什么？”

我不知道，这是不是巧合，但是我宁愿相信这两件事是有联系的，我得到了这么大的回报，不仅是因为我过去给予过朋友资助，也因为我发自内心地愿意给予。不是出于义务，也不是因为恐惧害怕，只是一种发自内心地想要给予那个小男孩一份礼物，我觉得我遇到了一个年轻而高贵的灵魂。

我可以诚实地告诉你，我的一生有过很多非常难熬的日子，既有经济上的磨难，也有情绪上的磨难，我们都曾有过这样难熬的日子，但是我再也没有感受过那种短缺的感觉，我以后也再不会有了。

本书的终极信息非常简单，就是我把车停在高速公路边上写在日记上的那一句话：财富的终极秘密是给予。

自由地给予，坦诚地给予，轻松地给予，享受地给予。甚至是在你觉得你没有任何东西可以给予的时候，也给予。你会发现，有一个像大海一样丰富的宝藏，就在你的心里，就在你的周围。生活总是会为你发生，而不是发生在你头上。欣赏上天赐予你的那份礼物吧，你富有了，从现在直到永远。

理解了这个真理，把我带回生命之源：上天造出我这个人是为了什么？上天选出我们所有人是为了什么？只有一个目的，成为一股行善的力量。明白这个真理，把我带回人生的最深刻意义，我一直不停地努力实现我的祈祷：所有我遇到的人，所有我有幸建立联系的人，每一天我都要成为上天赐给他们的福气。

尽管我可能没有跟你见过面，但我写本书也出于相同的想为你带来福气的状态，我在问，我在祈祷，每一章、每一页、每一个概念都能更好地帮助你更加深入地体验你是谁，更加深入地体验你这一生能够创造什么，能够给予什么。

我发自内心的愿望也是本书的目的，就是给你一种方法，提高和深化你的生活品质，以及所有那些你有福气去爱、去感动的人的生活品质。在这里，我能够为你提供服务，我感到非常荣幸。

我期待有一天能与你面对面地交流，也许是我在世界各地组织的某一场讲座活动上，或者只是跟你在街头偶然相遇。我会非常激动地听你讲，你是如何用这些原则让你的生活更幸福、更快乐的。

好了，我们就要分别了，我要深深地祝福你，祝愿你的生活永远富裕安康。祝愿你的人生充满快乐，充满激情，充满挑战，充满机会，充满成长，充满给予。我祝愿你的人生成就非凡。

我永远爱你，永远祝福你！

托尼·罗宾斯

充分感受生活，就在此时此地。用心体验每一件事，照顾好你自己和你的朋友。找点儿乐趣，要疯狂，要怪异。走出去，把事情搞砸！反正你肯定会出错的，所以你最好享受整个过程，不要太在意结果。抓住机会从错误中学习，找到问题所在，然后解决。不要想完美，只要成为一个优秀的个体就够了。

托尼·罗宾斯
